% #########################################################################
% #                       EXPLORACIÓN Y ANÁLISIS DE DATOS                 #
% #########################################################################

% Definición de nombres
\newcommand{\estudiante}{García Justel, Alan}
\newcommand{\titulo}{MÁSTER EN INGENIERÍA COMPUTACIONAL Y SISTEMAS INTELIGENTES}
\newcommand{\asignatura}{INTRODUCTION TO TIME SERIES DATA}
\newcommand{\portada}{common/no_signal.png}
\newcommand{\colorportada}{title_green}
\newcommand{\curso}{2024-2025}


% Notebook
\begin{document}
\newgeometry{bottom=2cm}

\begin{titlepage}
    % Logo de la universidad
    \begin{textblock*}{\textwidth}(2cm,0.4cm)
        \begin{center}
            \begin{minipage}{0.45\textwidth}
                \centering
                \includegraphics[width=\textwidth]{common/Logo_EHU.jpg}
            \end{minipage}\hfill
            \begin{minipage}{0.45\textwidth}
                \centering
                \includegraphics[width=0.8\textwidth]{common/Logo_KISA.png} 
            \end{minipage}
        \end{center}
    \end{textblock*}
    
    % Franja de color
    \begin{tikzpicture}[remember picture, overlay]
        \fill[\colorportada] (current page.north west) ++ (0,-3.01cm) rectangle (\paperwidth,-3cm);
    \end{tikzpicture}
    
    \begin{textblock*}{\paperwidth}(\dimexpr\parindent+\oddsidemargin+3em\relax,3.5cm)
        \begin{minipage}{\dimexpr\linewidth-7.5cm\relax}
            \color{white}
            \noindent\rule{\linewidth}{0cm}
            \textsf{ {\large \titulo}}
            \newline
            \newline \newline
            \textsf{\textbf{ {\Huge APUNTES DE ASIGNATURA }}}
        \end{minipage}
    \end{textblock*}
    
    % Nombre asignatura
    \vspace*{3.5cm}
    \begin{minipage}{\linewidth}
        \setlength{\baselineskip}{1.7\baselineskip}
        \centering
        \textsf{ \textbf{ {\LARGE \asignatura }}}
    \end{minipage}

    % Foto de portada
    \vspace*{0.5cm}
    \begin{figure}[H]
        \centering
        \includegraphics[width=10cm, height=8cm]{\portada}
    \end{figure}

    
    
    % Estudiante
    \vspace{0.2cm}
    \noindent {\footnotesize \textbf{Estudiante:} \estudiante}
    \newline
    \noindent\makebox[\linewidth]{\rule{\textwidth}{0.4pt}} % Línea horizontal

    % Curso y Fecha
    \vspace{0.1cm}    
    \noindent {\footnotesize \textbf{Curso: } \curso \hfill \textbf{Fecha:} \today }
\end{titlepage}

\restoregeometry
\setcounter{figure}{0} % Incluimos el título
\newpage
% Índices
\tableofcontents\thispagestyle{empty} \newpage
% \listoffigures\thispagestyle{empty}   %\newpage
% \listoftables\thispagestyle{empty}    %\newpage

\section{Introduction}
Time series data is information that represents one or more attributes (variables) over time (or another ordered dimension).
There are multiple aproximation for using a time series data for different tasks:
- Use the data directly or with a preprocessing step for noise and dimensionality reduction, split into windows of data, selecting specific parameter for represanting the time series data or other more sophisticated techniques like fourier methods.
- Use algorithms for specific data (convolutional networks for images, RNN for audio...).

\section{Time Series Forecasting}
The objective is to predict or to infer a future timestamp state for the time series.

$$
y_{t + a} = f(y_t, y_{t-1}, y_{t-2}, y_{t-3}, \dots, \text{error})
$$

Where $(y_{t}, y_{t}, y_{t})$ are the previous observations of the variables we want to predict (historical values). 
Also, it is possible to have other extra information ($(x_{1}, x_{2}, x_{3})$) that is useful for the forecasting:

$$
y_{t + a} = f(y_t, y_{t-1}, y_{t-2}, y_{t-3}, \dots, x_{1}, x_{2}, x_{3}, \dots, \text{error})
$$

\subsection{Classical lineal models}
The Classical lineal models has these assumptions:

\begin{itemize}
    \item \textbf{stochastic process}: an ordered sequence of random variables, where the
    order is typically given by time: $(Y_1, Y_2, Y_3, \dots, Y_n, \dots)$.
    \item \textbf{Time series}: a finite set of observations (a sample) from a real-valued
    stochastic process (univariate or multivariate).
\end{itemize}

We need a time series to be stationary in order to apply forecasting methods. A time series is stationary if these conditions are fulfilled:
$$
F(X_1, X_2, X_3, \dots, X_n) = P(X_1 < x_1, X_2 < x_2, X_3 < x_3, \dots, X_n < x_n)
$$
\begin{enumerate}
    \item First order stationary: $F(X_t) = F(X_{t +k})$. If the probability distribution of the first variable is the same as the second one, the same as the third one and so on.
    \item Second order stationary: $F(X_{t1}, X_{t2}) = F(X_{t1 + k}, X_{t2 + k})$.
\end{enumerate}

However, it is very difficult to have a full stationary time series, so a \textit{weak stationary} is defined:
\begin{enumerate}
    \item The mean is constant.
    \item The variance is constatnt.
    \item The covariance si constant in terms of a parameter ($\text{Cov}(X_i, X_{i+k})=f(k)$). It depends on time difference between the random variables.
\end{enumerate}

\textbf{How do we check if a time series is stationary?}
\begin{enumerate}
    \item Visual inspection
    \item Statistical tests
\end{enumerate}


\end{document}